%% Chapter3_SocietalImpact.tex
\chapter{Societal and Economic Impact}

\section{Societal & Consumer Protection Impact}
The global trade in counterfeit apparel disproportionately affects everyday consumers, online buyers, and secondary market sellers who fall victim to deceptive listing practices. Fake luxury apparel often involves hazardous chemical dyes, illegal synthetic materials, and unethical labor conditions \cite{oecd2020trade}. 

By deploying **VeriStyle**, consumers gain access to an accessible, real-time authenticity verification tool that democratizes expert-level apparel inspection. This protects consumers from financial exploitation, empowers honest resale sellers, and fosters a trustworthy online marketplace.

\section{Economic & Industry Impact}
Counterfeiting costs legitimate fashion brands, independent designers, and authorized e-commerce retailers hundreds of billions of dollars annually in lost revenue, brand equity erosion, and customer dispute resolution expenses. 

VeriStyle provides e-commerce platforms and resale marketplaces (e.g., eBay, Vestiaire Collective, Poshmark) with an automated, scalable REST API (`POST /api/analyze-authenticity`) capable of pre-screening thousands of listings per minute. By reducing fraudulent sales, businesses minimize credit card chargebacks, return handling costs, and legal liabilities associated with trademark infringement.

\section{Environmental & Sustainability Aspects}
The proliferation of low-quality counterfeit garments contributes significantly to fast-fashion waste, landfill pollution, and resource depletion. In contrast, authentic luxury apparel is designed for longevity, circular economy reuse, and sustained value retention. By safeguarding the authentic luxury resale ecosystem, VeriStyle encourages sustainable consumption patterns and extends the lifecycle of fashion goods.
